%==================================================================
% PERNYATAAN PENGGUNAAN KECERDASAN BUATAN DALAM PENULISAN
%
% Mengikuti kelaziman jurnal internasional (mis. Elsevier), pernyataan
% penggunaan AI pada proses PENULISAN diletakkan pada bagian tersendiri
% tepat SEBELUM Daftar Pustaka. Sebutkan nama alat, tujuan penggunaan,
% dan pernyataan tanggung jawab penulis.
%
% Penggunaan AI dalam proses PENELITIAN dijelaskan terpisah pada Bab 3
% (lihat \cref{sec:ai-penelitian}).
%
% Bila AI tidak dipakai untuk penulisan, beri tanda % pada baris
% %==================================================================
% PERNYATAAN PENGGUNAAN KECERDASAN BUATAN DALAM PENULISAN
%
% Mengikuti kelaziman jurnal internasional (mis. Elsevier), pernyataan
% penggunaan AI pada proses PENULISAN diletakkan pada bagian tersendiri
% tepat SEBELUM Daftar Pustaka. Sebutkan nama alat, tujuan penggunaan,
% dan pernyataan tanggung jawab penulis.
%
% Penggunaan AI dalam proses PENELITIAN dijelaskan terpisah pada Bab 3
% (lihat \cref{sec:ai-penelitian}).
%
% Bila AI tidak dipakai untuk penulisan, beri tanda % pada baris
% %==================================================================
% PERNYATAAN PENGGUNAAN KECERDASAN BUATAN DALAM PENULISAN
%
% Mengikuti kelaziman jurnal internasional (mis. Elsevier), pernyataan
% penggunaan AI pada proses PENULISAN diletakkan pada bagian tersendiri
% tepat SEBELUM Daftar Pustaka. Sebutkan nama alat, tujuan penggunaan,
% dan pernyataan tanggung jawab penulis.
%
% Penggunaan AI dalam proses PENELITIAN dijelaskan terpisah pada Bab 3
% (lihat \cref{sec:ai-penelitian}).
%
% Bila AI tidak dipakai untuk penulisan, beri tanda % pada baris
% %==================================================================
% PERNYATAAN PENGGUNAAN KECERDASAN BUATAN DALAM PENULISAN
%
% Mengikuti kelaziman jurnal internasional (mis. Elsevier), pernyataan
% penggunaan AI pada proses PENULISAN diletakkan pada bagian tersendiri
% tepat SEBELUM Daftar Pustaka. Sebutkan nama alat, tujuan penggunaan,
% dan pernyataan tanggung jawab penulis.
%
% Penggunaan AI dalam proses PENELITIAN dijelaskan terpisah pada Bab 3
% (lihat \cref{sec:ai-penelitian}).
%
% Bila AI tidak dipakai untuk penulisan, beri tanda % pada baris
% \include{xx-deklarasi-ai} di laporan.tex.
%==================================================================
\clearpage
\phantomsection
\addcontentsline{toc}{chapter}{PERNYATAAN PENGGUNAAN KECERDASAN BUATAN}
\begin{center}
    \textbf{\large PERNYATAAN PENGGUNAAN KECERDASAN BUATAN}\\
    \textbf{\large DALAM PROSES PENULISAN}\\[2em]
\end{center}

Selama proses penyusunan naskah Tugas Akhir ini, penulis menggunakan
\belumdiisi{sebutkan nama dan versi alat AI, mis. ChatGPT (GPT-4), Claude, atau
Grammarly} untuk \belumdiisi{sebutkan tujuannya, mis. membantu merapikan dan
memperbaiki keterbacaan bahasa, menata struktur penulisan, serta menerjemahkan
abstrak ke dalam Bahasa Inggris}. Setelah menggunakan alat tersebut, penulis
telah meninjau dan menyunting seluruh isinya sesuai kebutuhan dan bertanggung
jawab penuh atas isi naskah Tugas Akhir ini.

\vspace{0.8em}
\noindent Pernyataan ini tidak berlaku untuk penggunaan alat bantu dasar seperti
pemeriksa tata bahasa, ejaan, dan tanda baca. Kecerdasan buatan tidak dicantumkan
sebagai penulis maupun sebagai sumber rujukan dalam Tugas Akhir ini.

\vspace{2em}
\begin{flushright}
    \kota, \tglpernyataan\\[0.8em]
    Penulis,\\[2cm]
    \penulis\\
    NIM. \nim
\end{flushright}
\clearpage
 di laporan.tex.
%==================================================================
\clearpage
\phantomsection
\addcontentsline{toc}{chapter}{PERNYATAAN PENGGUNAAN KECERDASAN BUATAN}
\begin{center}
    \textbf{\large PERNYATAAN PENGGUNAAN KECERDASAN BUATAN}\\
    \textbf{\large DALAM PROSES PENULISAN}\\[2em]
\end{center}

Selama proses penyusunan naskah Tugas Akhir ini, penulis menggunakan
\belumdiisi{sebutkan nama dan versi alat AI, mis. ChatGPT (GPT-4), Claude, atau
Grammarly} untuk \belumdiisi{sebutkan tujuannya, mis. membantu merapikan dan
memperbaiki keterbacaan bahasa, menata struktur penulisan, serta menerjemahkan
abstrak ke dalam Bahasa Inggris}. Setelah menggunakan alat tersebut, penulis
telah meninjau dan menyunting seluruh isinya sesuai kebutuhan dan bertanggung
jawab penuh atas isi naskah Tugas Akhir ini.

\vspace{0.8em}
\noindent Pernyataan ini tidak berlaku untuk penggunaan alat bantu dasar seperti
pemeriksa tata bahasa, ejaan, dan tanda baca. Kecerdasan buatan tidak dicantumkan
sebagai penulis maupun sebagai sumber rujukan dalam Tugas Akhir ini.

\vspace{2em}
\begin{flushright}
    \kota, \tglpernyataan\\[0.8em]
    Penulis,\\[2cm]
    \penulis\\
    NIM. \nim
\end{flushright}
\clearpage
 di laporan.tex.
%==================================================================
\clearpage
\phantomsection
\addcontentsline{toc}{chapter}{PERNYATAAN PENGGUNAAN KECERDASAN BUATAN}
\begin{center}
    \textbf{\large PERNYATAAN PENGGUNAAN KECERDASAN BUATAN}\\
    \textbf{\large DALAM PROSES PENULISAN}\\[2em]
\end{center}

Selama proses penyusunan naskah Tugas Akhir ini, penulis menggunakan
\belumdiisi{sebutkan nama dan versi alat AI, mis. ChatGPT (GPT-4), Claude, atau
Grammarly} untuk \belumdiisi{sebutkan tujuannya, mis. membantu merapikan dan
memperbaiki keterbacaan bahasa, menata struktur penulisan, serta menerjemahkan
abstrak ke dalam Bahasa Inggris}. Setelah menggunakan alat tersebut, penulis
telah meninjau dan menyunting seluruh isinya sesuai kebutuhan dan bertanggung
jawab penuh atas isi naskah Tugas Akhir ini.

\vspace{0.8em}
\noindent Pernyataan ini tidak berlaku untuk penggunaan alat bantu dasar seperti
pemeriksa tata bahasa, ejaan, dan tanda baca. Kecerdasan buatan tidak dicantumkan
sebagai penulis maupun sebagai sumber rujukan dalam Tugas Akhir ini.

\vspace{2em}
\begin{flushright}
    \kota, \tglpernyataan\\[0.8em]
    Penulis,\\[2cm]
    \penulis\\
    NIM. \nim
\end{flushright}
\clearpage
 di laporan.tex.
%==================================================================
\clearpage
\phantomsection
\addcontentsline{toc}{chapter}{PERNYATAAN PENGGUNAAN KECERDASAN BUATAN}
\begin{center}
    \textbf{\large PERNYATAAN PENGGUNAAN KECERDASAN BUATAN}\\
    \textbf{\large DALAM PROSES PENULISAN}\\[2em]
\end{center}

Selama proses penyusunan naskah Tugas Akhir ini, penulis menggunakan
\belumdiisi{sebutkan nama dan versi alat AI, mis. ChatGPT (GPT-4), Claude, atau
Grammarly} untuk \belumdiisi{sebutkan tujuannya, mis. membantu merapikan dan
memperbaiki keterbacaan bahasa, menata struktur penulisan, serta menerjemahkan
abstrak ke dalam Bahasa Inggris}. Setelah menggunakan alat tersebut, penulis
telah meninjau dan menyunting seluruh isinya sesuai kebutuhan dan bertanggung
jawab penuh atas isi naskah Tugas Akhir ini.

\vspace{0.8em}
\noindent Pernyataan ini tidak berlaku untuk penggunaan alat bantu dasar seperti
pemeriksa tata bahasa, ejaan, dan tanda baca. Kecerdasan buatan tidak dicantumkan
sebagai penulis maupun sebagai sumber rujukan dalam Tugas Akhir ini.

\vspace{2em}
\begin{flushright}
    \kota, \tglpernyataan\\[0.8em]
    Penulis,\\[2cm]
    \penulis\\
    NIM. \nim
\end{flushright}
\clearpage
